
\section{Leverage Score Sampling\label{sec:Leverage-Score-Sampling}}

Here we give a way to reduce solving a tall linear regression system
into a sequence of submatrices:
\begin{thm}[\cite{cohen2015uniform}]
Given $A\in\R^{n\times d}$, let $T(m)$ be the cost of solving $B^{\top}Bx=c$,
for arbitrary $c\in\R^{d}$, over all submatrices $B$ consisting
of $m$ distinct rows of $A$. Then, we have that
\[
T(n)=O^{*}(\nnz(A)+T(O(d\log d))).
\]
\end{thm}

\begin{rem*}
Here we use $O^{*}$ to emphasize there is some dependence on $\log(\frac{1}{\epsilon})$
suppressed for notational simplicity. Lemma \ref{lem:Richardson}
shows that once we can solve the system with constant approximation,
we can repeat $\log(\frac{1}{\epsilon})$ times to get an $\epsilon$-accurate
solution.
\end{rem*}

\subsection{Leverage scores}

The key concept in this reduction is the leverage score.
\begin{defn}
Given a matrix $A\in\R^{n\times d}$, let $a^{\top}_{i}$ be the $i^{th}$
row of $A$. The leverage score of the $i^{th}$ row of $A$ is
\[
\sigma_{i}(A)\defeq a^{\top}_{i}(A^{\top}A)^{+}a_{i}.
\]
\end{defn}

Note that $\sigma(A)$ is the diagonal of the projection matrix $A(A^{\top}A)^{+}A^{\top}$.
Since $0\preceq A(A^{\top}A)^{+}A^{\top}\preceq I$, we have that
$0\leq\sigma_{i}(A)\leq1$. Moreover, since $A(A^{\top}A)^{+}A^{\top}$
is a projection matrix, the sum of $A$'s leverage scores (its trace)
is equal to the rank of $A$: 
\begin{align}
\sum^{n}_{i=1}\sigma_{i}(A)=\tr(A(A^{\top}A)^{+}A^{\top}) & =\text{rank}(A(A^{\top}A)^{+}A^{\top})=\text{rank}(A)\le d.\label{rank_bound}
\end{align}
The leverage score measures the ``importance'' of a row in forming
the row space of $A$. If a row has a component orthogonal to all
other rows, its leverage score is $1$. Removing it would decrease
the rank of $A$, completely changing its row space. The \emph{coherence}
of $A$ is $\norm{\sigma(A)}_{\infty}$. If $A$ has low coherence,
no particular row is especially important. If $A$ has high coherence,
it contains at least one row whose removal would significantly affect
the composition of $A$'s row space. The following two characterizations
help with this intuition. 
\begin{lem}
\label{leverage_score_opt} For all $A\in\R^{n\times d}$ and $i\in[n]$
we have that
\[
\sigma_{i}(A)=\min_{A^{\top}x=a_{i}}\norm x^{2}_{2}.
\]
where $a^{\top}_{i}$ is the $i^{th}$ row of $A$. 
\end{lem}

\begin{lem}
\label{leverage_score_opt_2} For all $A\in\R^{n\times d}$ and $i\in[n]$,
we have that $\sigma_{i}(A)$ is the smallest $t$ such that 
\begin{equation}
a_{i}a^{\top}_{i}\preceq t\cdot A^{\top}A.\label{eq:leverage_score_spectral_upper}
\end{equation}
\end{lem}

Sampling rows from $A$ according to their exact leverage scores gives
a spectral approximation for $A$ with high probability. Sampling
with leverage-score overestimates also suffices.
\begin{lem}
\label{lem:leveragescoresampling}Given a vector $u$ of leverage-score
overestimates, i.e., $\sigma_{i}(A)\le u_{i}$ for all $i$, define
\[
X=\frac{1}{p_{i}}a_{i}a^{\top}_{i}\quad\mbox{ with probability \quad\ }p_{i}=\frac{u_{i}}{\norm u_{1}}.
\]
For $T=\Omega(\frac{\norm u_{1}\log n}{\varepsilon^{2}})$, with probability
$1-\frac{1}{n^{O(1)}}$, we have that
\[
(1-\varepsilon)A^{\top}A\preceq\frac{1}{T}\sum^{T}_{i=1}X_{i}\preceq(1+\varepsilon)A^{\top}A
\]
where $X_{i}$ are independent copies of $X$.
\end{lem}

\begin{proof}
Note that $\E X=A^{\top}A$ and that 
\[
0\preceq X=\frac{1}{p_{i}}a_{i}a^{\top}_{i}\preceq\frac{\norm u_{1}}{\sigma_{i}}a_{i}a^{\top}_{i}\preceq\norm u_{1}\cdot A^{\top}A.
\]
Now, the statement simply follows from the matrix Chernoff bound with
$M_{k}=(A^{\top}A)^{-\frac{1}{2}}X_{k}(A^{\top}A)^{-\frac{1}{2}}$
and $R=\norm u_{1}$.
\end{proof}
Combining Lemma \ref{lem:leveragescoresampling} and Lemma \ref{lem:Richardson},
we have that
\begin{equation}
T(n)=\text{cost of computing all leverage scores }\sigma+O^{*}(\nnz(A)+T(d\log d))\label{eq:TmnS1}
\end{equation}
where we used that $\norm{\sigma}_{1}=O(d)$. However, computing $\sigma$
exactly is too expensive for many purposes. In \cite{spielman2011graph},
the authors showed that we can compute leverage scores approximately
by solving only polylogarithmically many regression problems. This
result uses the fact that 
\[
\sigma_{i}(A)=\norm{A(A^{\top}A)^{+}A^{\top}e_{i}}^{2}_{2}
\]
and that by the Johnson-Lindenstrauss Lemma these lengths are preserved
up to a multiplicative error if we project these vectors to a random
low-dimensional subspace.

In particular, this lemma shows that we can approximate $\sigma_{i}(A)$
via $\norm{\Pi A(A^{\top}A)^{+}A^{\top}e_{i}}^{2}_{2}$. The benefit
of this is that we can compute $\Pi A(A^{\top}A)^{+}$ by solving
$\frac{\log n}{\varepsilon^{2}}$ many linear systems. In other words,
we have that 
\[
\text{cost of approximating }\sigma_{i}=O^{*}(\nnz(A))+O^{*}(T(n)).
\]
Putting it into (\ref{eq:TmnS1}), we have this non-informative recurrence
\[
T(n)=O^{*}(\nnz(A)+T(n)+T(d\log d)).
\]
This is a chicken-and-egg problem. To solve an overdetermined system
faster, we want to use leverage scores to sample the rows. To approximate
the leverage scores, we need to solve the original overdetermined
system. (Even worse, we need to solve it a few times.)

\subsection{Uniform Sampling}

The key idea to break this chicken-and-egg problem is to use uniform
sampling of the rows of $A$. We define
\[
\sigma_{i,S}=a^{\top}_{i}(A^{\top}_{S}A_{S})^{-1}a_{i}
\]
where $A_{S}$ is $A$ restricted to rows in $S$. The set $S$ will
be a random sample of $k$ rows of $A$. Note that since $A^{\top}_{S}A_{S}\preceq A^{\top}A$,
the quantity $\sigma_{i,S}$ is an overestimate of $\sigma_{i}(A)$.
Hence, it suffices to bound $\sum_{i}\sigma_{i,S}$. The key lemma
is the following:
\begin{lem}
\label{lem:sigmaSi}We have that
\[
\E_{|S|=k}\sum^{n}_{i=1}\sigma_{i,S\cup\{i\}}\leq\frac{nd}{k}.
\]
\end{lem}

\begin{proof}
Note that
\begin{align*}
\E_{|S|=k}\sum^{n}_{i=1}\sigma_{i,S\cup\{i\}} & =\E_{|S|=k}\sum_{i\notin S}\sigma_{i,S\cup\{i\}}+\E_{|S|=k}\sum_{i\in S}\sigma_{i,S\cup\{i\}}.
\end{align*}
Note that $\sum_{i\in S}\sigma_{i,S\cup\{i\}}=\sum_{i\in S}\sigma_{i,S}\leq d$.
Hence, the second term is bounded by $d$. 

For the first term, we note that ``sample a set $S$ of size $k$,
then sample $i\notin S$'' is the same as ``sample a set $T$ of
size $k+1$, then sample $i\in T$''. Hence, we have
\begin{align*}
\E_{|S|=k}\E_{i\notin S}\sigma_{i,S\cup\{i\}} & =\E_{|T|=k+1}\E_{i\in T}\sigma_{i,T}\\
 & \leq\E_{|T|=k+1}\frac{d}{k+1}=\frac{d}{k+1}
\end{align*}

Hence, we have that
\[
\E_{|S|=k}\sum^{n}_{i=1}\sigma_{i,S\cup\{i\}}\leq\frac{d}{k+1}(n-k)+d=d\cdot\frac{n+1}{k+1}.
\]
\end{proof}
Next, using the Sherman-Morrison formula, we have that
\[
\sigma_{i,S\cup\{i\}}=\begin{cases}
\sigma_{i,S} & \text{if }i\in S\\
\frac{1}{1+\frac{1}{\sigma_{i,S}}} & \mbox{otherwise}
\end{cases}.
\]
Namely, we can compute $\sigma_{i,S\cup\{i\}}$ using $\sigma_{i,S}$.
Therefore, we have that
\[
\text{cost of approximating }\sigma_{i,S\cup\{i\}}=O^{*}(\nnz(A))+O^{*}(T(k)).
\]
Using Lemma \ref{lem:sigmaSi}, we have now
\[
T(n)=O^{*}(\nnz(A)+T(k)+T(\frac{nd}{k}\log d)).
\]
Picking $k=\sqrt{nd\log d}$, we have that
\[
T(n)=O^{*}(\nnz(A)+T(\sqrt{nd\log d})).
\]
Repeating this process $\tilde{O}(1)$ times, we have
\[
T(n)=O^{*}(\nnz(A)+T(O(d\log d))).
\]


